\section{协议筛选与残差时间：折叠刚性与类内区别的解析深度}\label{sec:fold-residual-time}

本节在不改变全文既有对象层级与证明链条的前提下，对两条已出现但分散于不同章节的接口作一次统一整理：其一是“设计空间—协议约束—刚性筛选”的筛选链，用以精确说明本文所用折叠并非任意选择，而是与二元读出、唯一可寻址、跨分辨率相容与对称性审计等接口条件共同锁定的最简可审计结构；其二是“折叠商结构—纤维残差—时间尺度”的解析链，用以把“类内区别”提升为可度量的残差预算，并将在线延迟、同步时间尾与高阶碰撞矩等对象统一解释为残差在时间方向的展开证书。\par
本节仅作系统性编排与定理链导航：所有数学主张均已在前文给出编号结论或标准引理；本节的表述不引入新的基本假设，亦不重复原有证明。

\subsection{设计空间与协议筛选：折叠为何被剩下}\label{subsec:protocol-screening-fold-survivor}

本小节把第 \ref{sec:spg}--\ref{sec:group-unification} 节中已出现的“二元读出—规范形—跨分辨率相容—对称性审计”口径，统一整理为一条协议筛选链，并将其与一般 Ostrowski 设计空间（小节 \ref{subsec:ostrowski-generalized-fold}）对齐。其目的在于澄清：本文采用的 Zeckendorf 折叠并非从众多可能规则中任取一例，而是在本文固定的可审计接口条件下，由设计空间内生地退化得到的最简幸存结构。\par
为避免重复证明，本小节的陈述以引用为主；必要的推导仅限于由定义直接推出的参数退化关系。

\subsubsection{设计空间：一般 Ostrowski 计数与“取值 \texorpdfstring{$\to$}{->} 规范形 \texorpdfstring{$\to$}{->} 截断”折叠族}
在小节 \ref{subsec:ostrowski-generalized-fold} 中，本文已将二元 Zeckendorf 折叠推广到一般无理数 \(\alpha=[0;a_1,a_2,\dots]\) 所诱导的 Ostrowski 计数系统：对每个分辨率 \(m\)，其原始数字空间与合法稳定类型空间分别为
\[
\Omega_m^{(\alpha)}=\prod_{n=1}^{m}\{0,1,\dots,a_n\},\qquad
X_m^{(\alpha)}\subseteq \Omega_m^{(\alpha)}
\]
（定义 \ref{def:ostrowski-entropy-gap}），并由 Ostrowski 表示定理给出唯一规范形 \(Z_\alpha(N)\) 与合法集合 \(X_m^{(\alpha)}\)（定义 \ref{def:Xm-alpha}）。相应的折叠族
\[
\Fold_m^{(\alpha)}(\omega):=\pi_m^{(\alpha)}\!\bigl(Z_\alpha(N_\alpha(\omega))\bigr)
\qquad(\text{定义 \ref{def:fold-alpha}})
\]
体现了本文固定的构造范式：“先取值、再取规范形、最后截断”。该范式的直接后果是跨分辨率函子性（命题 \ref{prop:fold-alpha-basic}-(4) 及定理 \ref{thm:normalize-then-truncate-functorial}）：在一般 Ostrowski 场景中，“先规范形再截断”与截断投影严格可交换，交换图是定义性成立而非附加假设。\par
本文的主线折叠 \(\Fold_m:\Omega_m\to X_m\) 则是上述族在黄金端点 \(\alpha=\varphi^{-1}\) 的二元退化（命题 \ref{prop:Xm-alpha-cardinality}，并见定义 \ref{def:fold-word}）。

\subsubsection{协议条款：六类可审计接口条件}
将本文全稿反复使用的“协议”要求，抽象为一组仅约束\emph{允许的折叠接口}的条款。以下条款不依赖具体实现细节，只要求输出对象在有限窗口与可核验语义下可闭合：
\begin{enumerate}
  \item[\textbf{(P1)}] \textbf{二元读出：} 宏观读出字母表固定为 \(\{0,1\}\)（第 \ref{sec:spg} 节）；允许的折叠与中间门不得引入更大字母表。
  \item[\textbf{(P2)}] \textbf{唯一可寻址：} 宏观地址必须具备唯一指称：对每个分辨率 \(m\)，稳定类型层应允许用有限对象给出可核验的单值地址化（与“地址先于概念值/空值 \(\mathsf{NULL}\)”口径相容，见第 \ref{sec:recursive-addressing} 节与第 \ref{sec:pom} 节的类型化讨论）。
  \item[\textbf{(P3)}] \textbf{有限窗口完备：} 分辨率 \(m\) 的宏观对象应由前 \(m\) 位完备决定，且其可寻址值域在该尺度上无孔洞（需存在双射证书把 \(X_m\) 对齐到连续区间）。
  \item[\textbf{(P4)}] \textbf{跨分辨率相容：} 折叠与截断/降阶应在可审计意义下相容：同一对象的多尺度读出必须形成可交换的比较图，从而避免“同一对象在不同分辨率下不一致”。
  \item[\textbf{(P5)}] \textbf{最坏情形全定义：} 折叠与降阶规则必须对所有输入定义，并在最坏输入上保持协议条款的可核验性；不允许以“避开坏输入”维持一致性。
  \item[\textbf{(P6)}] \textbf{对称性审计：} 协议允许的对称必须可枚举并可核验；未经授权的对称会诱发地址冲突或证书不可判别，因而须被筛除到可审计的残余稳定子。
\end{enumerate}

\subsubsection{筛选链：二元读出迫使 Ostrowski 退化到黄金端点}
在 Ostrowski 设计空间内，(P1) 直接把连分数参数锁定到黄金端点。其理由仅依赖定义：由 \(\Omega_m^{(\alpha)}=\prod_{n=1}^{m}\{0,1,\dots,a_n\}\)（定义 \ref{def:fold-alpha}）可知，若对所有 \(m\) 都要求 \(\Omega_m^{(\alpha)}=\{0,1\}^m\)，则必有 \(a_n=1\) 对所有 \(n\) 成立，从而
\[
\alpha=[0;1,1,1,\dots]=\varphi^{-1}.
\]
在该退化下，合法集合 \(X_m^{(\alpha)}\) 亦退化为禁止相邻 \(11\) 的黄金均值语法 \(X_m\)（命题 \ref{prop:Xm-alpha-cardinality}；并见小节 \ref{subsec:folding-fibonacci-stable-syntax}）。因此，黄金端点并非“偏好选择”，而是二元字母表与 Ostrowski 合法性在本文范式中的共同交点。

\begin{remark}[二元读出不等于二进制权重]\label{rem:protocol-binary-not-base2}
(P1) 限制的是读出字母表，而非值权重。本文二元折叠以 Fibonacci 权重取值并取 Zeckendorf 规范形（定义 \ref{def:fold-word}），其稳定语法由禁词 \((11)\) 决定（小节 \ref{subsec:folding-fibonacci-stable-syntax}）。若改用“二进制权重 + 满移位”并令折叠退化为恒等，则纤维多重度恒为 \(1\)，从而无法产生本文所需的非平凡纤维统计、规范差审计与信息账本接口（参见小节 \ref{subsec:pom-information-ledger} 及小节 \ref{subsec:folding-stable-compression-resolution} 的分解口径）。
\end{remark}

\subsubsection{有限窗口完备：前缀双射与连续值域证书}
在黄金端点的二元语法下，(P3) 的“无孔洞完备性”可由前缀双射证书给出。引理 \ref{lem:zeckendorf-range-bijection} 证明：存在双射
\[
V_m:X_m\overset{\sim}{\longrightarrow}\{0,1,\dots,F_{m+2}-1\},
\]
从而 \(X_m\) 的地址值域为连续区间并由前 \(m\) 位完备决定（亦见第 \ref{sec:emergent-arithmetic} 节把该双射提升为分辨率模结构的值—地址同构）。该双射在本文语义中提供了可审计的“尺度 \(m\) 完备性”证书：它排除了任何必须依赖远端尾部才能决定前缀取值的规则。

\subsubsection{跨分辨率相容与最坏情形：从交换图到全局重写}
在一般 Ostrowski 折叠中，(P4) 的跨分辨率相容性是定义性成立（定理 \ref{thm:normalize-then-truncate-functorial}）。在二元 Zeckendorf 折叠中，朴素截断 \(r_{m+1\to m}\) 对应“先截断再规范化”的顺序，因此一般不交换；而折叠感知 restriction
\[
\widetilde r_{m+1\to m}:=\pi_{m+1\to m}\circ \Fold_{m+1}
\]
正是把降阶顺序纠正为“先规范形再截断”的定义性实现，并且在自然性约束下是唯一修正（定理 \ref{thm:normalize-then-truncate-functorial}；并见命题 \ref{prop:fold-omega-commute}）。\par
将 (P4) 与 (P5) 合并后，出现一条更强的结构后果：若仍试图用“仅修正末端常数位”的局部尾补丁来恢复交换，则在最坏输入上必失败。其可审计证据是规范差 \(G_m\) 的最坏端满长翻转（定理 \ref{thm:fold-gauge-anomaly-max}）及“尾补丁不完备”推论（推论 \ref{cor:fold-tail-patch-incomplete}）。因此，在本文固定的折叠范式下，“跨尺度一致性 \(\Rightarrow\) 全局重写/跨位传播”不是实现偏好，而是协议条款在最坏情形下的结构性强制。

\subsubsection{对称性审计：window-\texorpdfstring{$6$}{6} 锚点下的残余稳定子}
(P6) 在本文中以“可枚举的有限证书”形式落实：window-\(6\) 的折叠标签 \(F_6:\Omega_6\to X_6\) 可视为 \(6\) 维超立方体顶点上的标签函数，在几何自同构群 \(\mathrm{Aut}(Q_6)\) 作用下引入几何稳定子
\[
\mathrm{Stab}_{\mathrm{geo}}(F_6):=\{\sigma\in\mathrm{Aut}(Q_6):F_6\circ\sigma=F_6\}.
\]
定理 \ref{thm:foldbin6-geo-stabilizer-z2}（亦见推论 \ref{cor:terminal-foldbin6-geo-stabilizer-z2}）给出：该稳定子为阶 \(2\) 群，仅含恒等与唯一非平凡对合元；并且其在二进制编号上的可审计指纹为掩码位移 \(\pm 34\)（推论 \ref{cor:foldbin6-geo-mask-34}）。该结果将“允许的对称”压缩到可核验的 \(\ZZ_2\) 残余，从而与 (P2) 的唯一可寻址性兼容。

\paragraph{综合陈述（接口意义下的“幸存”）}
综上，在本文采用的设计空间（一般 Ostrowski 折叠族，定义 \ref{def:fold-alpha}）与协议条款 (P1)--(P6) 的共同约束下，二元 Zeckendorf 折叠可被视为一个由接口闭包强制选出的最简模型：二元读出把参数退化到黄金端点并锁定禁词 \((11)\) 的稳定语法（命题 \ref{prop:Xm-alpha-cardinality} 与小节 \ref{subsec:folding-fibonacci-stable-syntax}），前缀双射给出有限窗完备的连续值域证书（引理 \ref{lem:zeckendorf-range-bijection}），跨分辨率相容性迫使采用“先规范化再截断”的降阶并在最坏情形上强制全局重写（定理 \ref{thm:normalize-then-truncate-functorial}、定理 \ref{thm:fold-gauge-anomaly-max}），而对称性筛选在 window-\(6\) 锚点处将几何稳定子削减为可审计的唯一 \(\ZZ_2\) 残余（定理 \ref{thm:foldbin6-geo-stabilizer-z2}）。\par
在该意义下，“折叠为何被剩下”应被理解为接口闭合的刚性结果：若试图在保持上述协议接口的同时任意改造折叠机制，则必在至少一处破坏二元读出、唯一可寻址、有限窗完备、跨尺度相容、最坏情形全定义或对称性审计的可核验条件。


\subsection{残差展开与时间尺度：纤维自由度的解析深度}\label{subsec:residual-unfolding-time-depth}

本小节以折叠商结构为起点，把“类内区别”提升为可审计的残差量，并将在线延迟、同步时间尾与高阶碰撞矩等对象统一解释为残差在时间方向上的展开证书。其核心锚点为信息账本恒等式（定理 \ref{thm:pom-kl-ledger}）及其纤维 KL 分解（引理 \ref{lem:pom-fiberwise-kl}）。\par
本节中“时间”首先仍按第 \ref{sec:spg} 节的协议固定为单生成元作用的迭代计数；而“时间尺度/时间深度”指在折叠读出制度下，为解析纤维残差所必须支付的最小迭代深度（或等价的记忆深度/同步深度）。本小节不引入新的熵/距离公理；一切记号沿用第 \ref{sec:folding} 节与第 \ref{sec:pom} 节。

\subsubsection{折叠商结构与纤维分解}\label{subsubsec:frt-quotient-fibers}
固定分辨率 \(m\)。窗口微态空间为 \(\Omega_m=\{0,1\}^m\)，稳定类型空间为黄金均值语法块
\[
X_m:=\{w\in\{0,1\}^m:\text{不含子词 }11\}
\]
（小节 \ref{subsec:folding-fibonacci-stable-syntax}）。折叠算子 \(\Fold_m:\Omega_m\to X_m\) 由“Fibonacci 加权值 \(\to\) Zeckendorf 规范形 \(\to\) 截断”给出（定义 \ref{def:fold-word}），并满足规范性、幂等性与满射性（命题 \ref{prop:fold-basic}）。\par
由此诱导等价关系
\[
\omega\sim_m\omega'\quad\Longleftrightarrow\quad \Fold_m(\omega)=\Fold_m(\omega'),
\]
其等价类（纤维）即 \(\Fold_m^{-1}(x)\ (x\in X_m)\)，并以纤维多重度
\[
d_m(x):=\bigl|\Fold_m^{-1}(x)\bigr|
\]
定量刻画类内自由度（小节 \ref{subsec:folding-stable-compression-resolution}）。在均匀微态基线下，折叠推前分布为
\[
w_m(x)=\frac{d_m(x)}{2^m},
\]
并由压缩率 \(2^m/|X_m|=2^m/F_{m+2}\) 表明：当 \(m\) 增大时折叠必然为多对一，纤维残差在尺度上不可忽略（同上小节）。

\subsubsection{残差：信息账本的两种等价读法}\label{subsubsec:frt-residual-ledger}
信息账本在第 \ref{sec:pom} 节已被固定为一条严格恒等式。为便于本节的语义对齐，回顾其接口记号（小节 \ref{subsec:pom-information-ledger}）：对任意微观分布 \(\mu\in\Delta(\Omega_m)\)，令其稳定推前为 \(\pi:=P_m\mu=(\Fold_m)_*\mu\in\Delta(X_m)\)，并令 \(\mu^e:=Q_m\pi\) 为“在每条纤维上条件均匀”的提升。则定理 \ref{thm:pom-kl-ledger} 给出
\[
\boxed{\;
H(\mu)=H(\pi)+\mathbb{E}_{\pi}\!\bigl[\log d_m(X)\bigr]\;-\;D_{\mathrm{KL}}(\mu\|\mu^e).
\;}
\]
其中（i）\(H(\pi)\) 为稳定类型的可见熵；（ii）\(\mathbb{E}_{\pi}[\log d_m(X)]\) 为纤维规模预算（不可见容量）；（iii）\(D_{\mathrm{KL}}(\mu\|\mu^e)\) 为同一宏观下微观分布偏离“纤维均匀基线”的可审计结构残差。引理 \ref{lem:pom-fiberwise-kl} 进一步将该 KL 项分解为纤维 KL 的平均：
\[
D_{\mathrm{KL}}(\mu\|\mu^e)=\sum_{x\in X_m}\pi(x)\,D_{\mathrm{KL}}(\mu_x\|u_x),
\]
其中 \(\mu_x\) 是纤维条件分布，\(u_x\) 为纤维上的均匀分布。\par
因此，“残差”在本文口径下天然具有两种互补读法：一方面，条件熵差
\[
H(\mu)-H(\pi)
\]
刻画宏观读出无法解释的类内不确定性；另一方面，纤维预算与纤维 KL 给出其精确分解
\[
\boxed{\;
H(\mu)-H(\pi)=\mathbb{E}_{\pi}\!\bigl[\log d_m(X)\bigr]\;-\;D_{\mathrm{KL}}(\mu\|\mu^e).
\;}
\]
在均匀微态基线（\(\mu\) 在每条纤维上条件均匀）下有 \(D_{\mathrm{KL}}(\mu\|\mu^e)=0\)，故残差完全由纤维预算 \(\mathbb{E}_{\pi}[\log d_m]\) 承载；而当纤维内部仍含结构时，平均纤维 KL 精确刻画“预算扣除量”。

\begin{proposition}[条件均匀提升的 KL 等距嵌入]\label{prop:frt-conditional-uniform-lift-kl-isometry}
\leanverified{paper\_frt\_conditional\_uniform\_lift\_kl\_isometry}
固定 \(m\ge 1\)，并沿用 \(Q_m:\Delta(X_m)\to\Delta(\Omega_m)\) 的条件均匀提升。
则对任意 \(\pi,\eta\in\Delta(X_m)\) 有严格恒等式
\[
\boxed{\;
D_{\mathrm{KL}}(Q_m\pi\ \|\ Q_m\eta)=D_{\mathrm{KL}}(\pi\ \|\ \eta)\; }.
\]
\leanverified{paper\_frt\_conditional\_uniform\_lift\_kl\_isometry}
\end{proposition}
\begin{proof}
按定义，对 \(a\in\Omega_m\) 写 \(x_a:=\Fold_m(a)\) 与 \(d(x):=d_m(x)\)，则
\(
(Q_m\pi)(a)=\pi(x_a)/d(x_a)
\)、
\(
(Q_m\eta)(a)=\eta(x_a)/d(x_a)
\)。
因此
\[
\begin{aligned}
D_{\mathrm{KL}}(Q_m\pi\|Q_m\eta)
&=\sum_{a\in\Omega_m}\frac{\pi(x_a)}{d(x_a)}\log\frac{\pi(x_a)/d(x_a)}{\eta(x_a)/d(x_a)}
=\sum_{a\in\Omega_m}\frac{\pi(x_a)}{d(x_a)}\log\frac{\pi(x_a)}{\eta(x_a)}\\
&=\sum_{x\in X_m}\left(\sum_{a:\,x_a=x}\frac{\pi(x)}{d(x)}\right)\log\frac{\pi(x)}{\eta(x)}
=\sum_{x\in X_m}\pi(x)\log\frac{\pi(x)}{\eta(x)}
=D_{\mathrm{KL}}(\pi\|\eta).
\end{aligned}
\]
证毕。
\end{proof}

\begin{proposition}[折叠 KL 几何的 Pythagoras 分解与 I-投影唯一性]\label{prop:frt-fold-kl-pythagoras-iprojection}
固定 \(m\ge 1\)，对任意 \(\mu\in\Delta(\Omega_m)\) 令 \(\pi:=P_m\mu\in\Delta(X_m)\)，并记条件均匀提升 \(\mu^e:=Q_m\pi\)。
则对任意 \(\eta\in\Delta(X_m)\) 有严格分解
\[
\boxed{\;
D_{\mathrm{KL}}(\mu\ \|\ Q_m\eta)
=D_{\mathrm{KL}}(\mu\ \|\ \mu^e)+D_{\mathrm{KL}}(\pi\ \|\ \eta)\; }.
\]
从而 \(\mu^e\) 是流形 \(\{Q_m\eta:\eta\in\Delta(X_m)\}\) 上对 \(\mu\) 的唯一 I-投影：它唯一最小化 \(D_{\mathrm{KL}}(\mu\|Q_m\eta)\)，并且最小值等于 \(D_{\mathrm{KL}}(\mu\|\mu^e)\)。
\leanverified{paper\_frt\_fold\_kl\_pythagoras\_iprojection}
\end{proposition}
\begin{proof}
仍写 \(x_a=\Fold_m(a)\)、\(d(x)=d_m(x)\)，则
\(
\mu^e(a)=\pi(x_a)/d(x_a)
\)、
\(
(Q_m\eta)(a)=\eta(x_a)/d(x_a)
\)，从而
\(
\log\frac{\mu^e(a)}{(Q_m\eta)(a)}=\log\frac{\pi(x_a)}{\eta(x_a)}
\)。
于是
\[
\begin{aligned}
D_{\mathrm{KL}}(\mu\|Q_m\eta)-D_{\mathrm{KL}}(\mu\|\mu^e)
&=\sum_{a\in\Omega_m}\mu(a)\log\frac{\mu^e(a)}{(Q_m\eta)(a)}\\
&=\sum_{a\in\Omega_m}\mu(a)\log\frac{\pi(x_a)}{\eta(x_a)}
=\sum_{x\in X_m}\left(\sum_{a:\,x_a=x}\mu(a)\right)\log\frac{\pi(x)}{\eta(x)}\\
&=\sum_{x\in X_m}\pi(x)\log\frac{\pi(x)}{\eta(x)}
=D_{\mathrm{KL}}(\pi\|\eta),
\end{aligned}
\]
即得所示分解。
由 \(D_{\mathrm{KL}}(\pi\|\eta)\ge 0\) 且等号当且仅当 \(\eta=\pi\) 可知最小化点唯一，并且 \(\arg\min_{\eta}D_{\mathrm{KL}}(\mu\|Q_m\eta)=\pi\)，对应 \(Q_m\pi=\mu^e\)。证毕。
\end{proof}

\subsubsection{时间作为解析深度：前缀超度量与分辨率分离}\label{subsubsec:frt-time-as-depth}
第 \ref{sec:spg} 节已将“时间”固定为迭代计数，并在符号柱集上引入前缀超度量 \(d_{\mathrm{pre}}\)（段落 \ref{par:spg-prefix-ultrametric}），其球等同于柱集，从而“前缀细化 \(\Leftrightarrow\) 半径收缩”。在地址语义侧，dyadic 前缀相容深度与地址超度量 \(d_{\mathrm{addr}}\) 的定义见定义 \ref{def:xi-dyadic-prefix-depth-ultrametric}。\par
本节的“时间尺度”即在上述前缀几何上引入折叠读出后的\emph{可分离层级}：纤维内差别并非在固定分辨率上一次性可见，而是随着分辨率 \(m\) 的推进逐层外化为可判别差别。为避免与正文中自同构 \(\tau\) 的记号冲突，定义如下分离深度。
\begin{definition}[折叠分离深度]\label{def:frt-fold-separation-depth}
对任意两条一侧微态流 \(\omega,\omega'\in\{0,1\}^{\mathbb{N}}\)，记其长度 \(m\) 前缀为 \(\omega|_m,\omega'|_m\in\Omega_m\)。定义折叠分离深度
\[
m_{\mathrm{sep}}(\omega,\omega')
:=\inf\{\,m\ge 1:\ \Fold_m(\omega|_m)\ne \Fold_m(\omega'|_m)\,\}\in\NN\cup\{\infty\}.
\]
\end{definition}

\begin{proposition}[投影相等的纤维 lift 审计]\label{prop:distill-euclid-elements-frt-projected-equality-lift-audit}
\textbf{Dependency status: rigidity.}
固定 \(m\ge1\)。对 \(\mu,\nu\in\Delta(\Omega_m)\)，令
\[
\pi:=P_m\mu,
\qquad
\eta:=P_m\nu.
\]
假设 \(\pi=\eta\)。对每个 \(x\in X_m\) 记纤维
\[
F_x:=\Fold_m^{-1}(x).
\]
当 \(\pi(x)>0\) 时，记 \(\mu_x,\nu_x\) 为 \(F_x\) 上的条件分布；当 \(\pi(x)=0\) 时，该纤维在 \(\mu\) 与 \(\nu\) 下均为零质量。则：
\begin{enumerate}
  \item
  \[
  \mu=\nu
  \quad\Longleftrightarrow\quad
  \mu_x=\nu_x\ \text{for every }x\in X_m\text{ with }\pi(x)>0;
  \]
  \item
  \[
  \mu=Q_m\pi
  \quad\Longleftrightarrow\quad
  D_{\mathrm{KL}}(\mu\|Q_m\pi)=0
  \quad\Longleftrightarrow\quad
  \mu_x=u_x\ \text{for every }x\text{ with }\pi(x)>0,
  \]
  where \(u_x\) is the uniform distribution on \(F_x\);
  \item consequently, projected equality \(P_m\mu=P_m\nu\) upgrades to the common lifted equality
  \[
  \mu=\nu=Q_m\pi
  \]
  if and only if both fibre residuals vanish:
  \[
  D_{\mathrm{KL}}(\mu\|Q_m\pi)=0,
  \qquad
  D_{\mathrm{KL}}(\nu\|Q_m\pi)=0.
  \]
\end{enumerate}
Thus equality after folding is only projected congruence until the fibre residual ledger is separately audited.
\end{proposition}

\begin{definition}[归一化折叠分离前沿]\label{def:distill-gromov-frt-normalized-separation-front}
令 \((\omega^{(n)},\omega'^{(n)})_{n\ge1}\) 为一列一侧微态流对，令 \(R_n\in\mathbb N\) 且 \(R_n\to\infty\)。记
\[
m_n:=m_{\mathrm{sep}}(\omega^{(n)},\omega'^{(n)})\in\mathbb N\cup\{\infty\},
\]
其中 \(m_{\mathrm{sep}}\) 如定义 \ref{def:frt-fold-separation-depth}。若 \(m_n<\infty\)，定义归一化分离前沿
\[
\Theta_n(t):=\mathbf 1\{\,m_n\le \lfloor tR_n\rfloor\,\},\qquad t\ge0.
\]
该函数记录在尺度 \(R_n\) 下，折叠读出首次分离是否已经发生。
\end{definition}

\begin{proposition}[有限测试失明与非坍缩分离前沿]\label{prop:distill-gromov-frt-normalized-separation-front}
沿用定义 \ref{def:distill-gromov-frt-normalized-separation-front}。假设每个 \(m_n<\infty\)，且存在 \(\lambda\in(0,\infty)\) 使
\[
\frac{m_n}{R_n}\longrightarrow \lambda.
\]
则成立：
\begin{enumerate}
  \item 对任意 \(t<\lambda\)，有 \(\Theta_n(t)=0\) 对所有充分大的 \(n\) 成立；对任意 \(t>\lambda\)，有 \(\Theta_n(t)=1\) 对所有充分大的 \(n\) 成立。
  \item 对任意固定深度 \(h\ge1\) 与任意有限读出 \(\tau_h:X_h\to S\)，若 \(n\) 充分大，则
  \[
  \tau_h\!\bigl(\Fold_h(\omega^{(n)}|_h)\bigr)
  =
  \tau_h\!\bigl(\Fold_h(\omega'^{(n)}|_h)\bigr).
  \]
\end{enumerate}
因此，该序列在每个固定有限分辨率测试下最终不可分，但在归一化尺度上保留非平凡分离前沿；这正是残差展开深度的粗尺度非坍缩见证。
\end{proposition}

\begin{theorem}[有界局部复杂度与无界残差深度的归一化锥骨架抽取]\label{distill:gromov-coarse_limit_obstruction_families-frt-normalized-residual-cone-extraction}
令 \(T_n\to\infty\) 为一列残差解析深度。对每个 \(n\)，令
\[
P_n:=\{0,1,\ldots,T_n\}
\]
并赋予归一化路径度量
\[
d_n(i,j):=\frac{|i-j|}{T_n}.
\]
设 \(\Sigma\) 为有限集合，并给定局部证书标签
\[
\ell_n:P_n\longrightarrow\Sigma.
\]
该有限标签假设表示残差展开在每个单位时间窗口内只有有界多种局部类型；若需要半径 \(r\) 的局部类型，只需把标签替换为有限块标签 \(\Sigma^{2r+1}\)。则存在子列 \(n_j\) 与映射
\[
\ell_\infty:\QQ\cap[0,1]\longrightarrow\Sigma
\]
使得对每个 \(t\in\QQ\cap[0,1]\)，序列
\[
\ell_{n_j}\bigl(\lfloor tT_{n_j}\rfloor\bigr)
\]
最终恒定，并且其最终值等于 \(\ell_\infty(t)\)。同时，度量空间 \((P_{n_j},d_{n_j},0,T_{n_j})\) 收敛到带两个标记点的区间
\[
([0,1],|\cdot|,0,1),
\]
且两标记点距离恒为
\[
|1-0|=1.
\]
因此，若 \(T_n\) 是某一折叠读出、同步读出或 residual lift 任务的最小解析深度，则任何固定深度 \(h\) 的有限测试在归一化尺度下坍缩到根点，因为
\[
\operatorname{diam}_{d_n}\{0,1,\ldots,h\}=\frac{h}{T_n}\longrightarrow0,
\]
而完整残差任务仍保留端点距离 \(1\)。这给出一个由本文自身归一化度量定义的非坍缩残差锥骨架。
\end{theorem}

\begin{proposition}[残差统计的尾纤维下降判别]\label{distill:cech_cohomology-refinement_stability_separation-frt-tail-fibre-descent-criterion}
固定 \(m\ge h\ge1\)，并令
\[
\rho_{m,h}:X_m\longrightarrow X_h,\qquad
\rho_{m,h}(x_1,\ldots,x_m):=(x_1,\ldots,x_h)
\]
为稳定代表上的 fold-aware 前缀限制。设 \(\varphi:X_m\to\RR\) 是一个有限分辨率统计量。以下三项等价：
\begin{enumerate}
  \item \(\varphi\) 是深度 \(h\) 的细化稳定统计量，即存在 \(\bar\varphi:X_h\to\RR\) 使
  \[
  \varphi=\bar\varphi\circ\rho_{m,h}.
  \]
  \item \(\varphi\) 在每条尾纤维上常值：若 \(\rho_{m,h}(x)=\rho_{m,h}(x')\)，则
  \[
  \varphi(x)=\varphi(x').
  \]
  \item 对任意概率测度 \(\mu,\nu\in\Delta(X_m)\)，若
  \[
  (\rho_{m,h})_\#\mu=(\rho_{m,h})_\#\nu,
  \]
  则
  \[
  \sum_{x\in X_m}\varphi(x)\mu(x)=
  \sum_{x\in X_m}\varphi(x)\nu(x).
  \]
\end{enumerate}
若这些条件失败，则存在一对
\[
(x,x')\in X_m\times X_m,\qquad \rho_{m,h}(x)=\rho_{m,h}(x'),\qquad \varphi(x)\ne\varphi(x'),
\]
并且两点测度 \(\delta_x,\delta_{x'}\) 在深度 \(h\) 下完全相同而 \(\varphi\)-统计不同。因此该失败是尾部 residual gap ledger，而不是深度 \(h\) 上的刚性不变量。
\end{proposition}
\begin{proof}
若 \(\varphi=\bar\varphi\circ\rho_{m,h}\)，则同一 \(
\rho_{m,h}\)-纤维上的两点显然有相同 \(\varphi\)-值，故第一项推出第二项。

若第二项成立，则可定义
\[
\bar\varphi(y):=\varphi(x)\qquad(x\in\rho_{m,h}^{-1}(y)).
\]
由于 \(\varphi\) 在每条纤维上常值，该定义与代表元无关。于是
\(\varphi=\bar\varphi\circ\rho_{m,h}\)，故第二项推出第一项。

第一项推出第三项如下。若 \(\varphi=\bar\varphi\circ\rho_{m,h}\)，则
\[
\sum_{x\in X_m}\varphi(x)
\mu(x)
=
\sum_{y\in X_h}\bar\varphi(y)(\rho_{m,h})_\#\mu(y),
\]
右端只依赖推前测度。因而推前相等的 \(\mu,\nu\) 给出相同期望。

最后证明第三项推出第二项。若第二项失败，则存在 \(x,x'\in X_m\) 满足
\(\rho_{m,h}(x)=\rho_{m,h}(x')\) 但 \(\varphi(x)\ne\varphi(x')\)。取
\(\mu=\delta_x\)、\(\nu=\delta_{x'}\)。二者在 \(X_h\) 上的推前都是同一个点质量，但 \(\varphi\)-期望分别为 \(\varphi(x)\) 与 \(\varphi(x')\)，矛盾于第三项。故第三项成立时第二项必成立。

失败情形正由上述构造给出。它说明深度 \(h\) 的折叠读出没有看见尾部差别，而深度 \(m\) 的统计量仍依赖该尾部 residual。因此该现象只能登记为有限分辨率尾纤维缺口；若要把它提升为稳定不变量，必须附加下降证书 \(\varphi=\bar\varphi\circ\rho_{m,h}\)。证毕。
\end{proof}

\begin{proof}
集合 \(\QQ\cap[0,1]\) 可数，取枚举 \(t_1,t_2,\ldots\)。由于 \(
\Sigma\) 有限，序列
\[
\ell_n\bigl(\lfloor t_1T_n\rfloor\bigr)
\]
存在常值子列。沿该子列再对 \(t_2\) 抽取常值子列，如此递归。取对角子列 \((n_j)\)。则对每个固定 \(t_a\)，当 \(j\) 充分大时，\(
\ell_{n_j}(\lfloor t_aT_{n_j}\rfloor)
\) 已经在第 \(a\) 步之后被固定，故最终恒定。定义 \(\ell_\infty(t_a)
\) 为该最终值，即得所需的有理点标签极限。

度量收敛是显式的。定义
\[
\iota_n:P_n\longrightarrow[0,1],\qquad \iota_n(i):=\frac{i}{T_n}.
\]
则 \(\iota_n\) 为等距嵌入，因为
\[
|\iota_n(i)-\iota_n(j)|=\frac{|i-j|}{T_n}=d_n(i,j).
\]
并且 \(\iota_n(P_n)\) 是 \([0,1]\) 中网格步长为 \(1/T_n\) 的有限子集，其 Hausdorff 距离至多 \(1/T_n\)，故趋于 \(0\)。根点 \(0\) 与端点 \(T_n\) 分别被送到 \(0\) 与 \(1\)。于是带标记点的度量空间收敛到 \(([0,1],|\cdot|,0,1)\)。

最后，对任意固定 \(h\)，集合 \(\{0,
1,\ldots,h\}\) 在 \(d_n\) 下的直径为 \(h/T_n\)，随 \(n\to\infty\) 趋于 \(0\)。因此任何只读取固定有限深度的审计，在该归一化尺度下只能看见根点的无穷小邻域；但残差任务的终端点始终与根点相距 \(1\)。若 \(T_n\) 是最小解析深度，则该端点正是任务首次可闭合或首次可分离的位置，故极限区间不是平凡点，而是记录了有限测试失明与全深度非坍缩之间的结构差。证毕。
\end{proof}

\begin{proof}
若 \(t<\lambda\)，取 \(\epsilon>0\) 使 \(t<\lambda-\epsilon\)。由 \(m_n/R_n\to\lambda\)，充分大时有
\[
\frac{m_n}{R_n}>\lambda-\epsilon>t,
\]
从而 \(m_n>tR_n\ge \lfloor tR_n\rfloor\)，故 \(\Theta_n(t)=0\)。若 \(t>\lambda\)，同理取 \(\epsilon>0\) 使 \(t>\lambda+\epsilon\)。充分大时
\[
\frac{m_n}{R_n}<\lambda+\epsilon<t,
\]
于是 \(m_n<tR_n\)，并因 \(m_n\) 为整数而有 \(m_n\le\lfloor tR_n\rfloor\) 对充分大 \(n\) 成立，故 \(\Theta_n(t)=1\)。

对第二条，固定 \(h\)。由于 \(m_n/R_n\to\lambda>0\) 且 \(R_n\to\infty\)，有 \(m_n\to\infty\)。因此充分大时 \(m_n>h\)。按定义 \ref{def:frt-fold-separation-depth}，在所有 \(1\le k<h+1\) 的深度上两条流尚未被折叠读出区分，特别地
\[
\Fold_h(\omega^{(n)}|_h)=\Fold_h(\omega'^{(n)}|_h).
\]
任意 \(\tau_h\) 作用后仍相等。证毕。
\end{proof}

\begin{proof}
For any \(a\in F_x\) with \(\pi(x)>0\), the disintegration formula gives
\[
\mu(a)=\pi(x)\mu_x(a),
\qquad
\nu(a)=\pi(x)\nu_x(a).
\]
If \(\pi(x)=0\), then both \(\mu\) and \(\nu\) assign zero mass to the whole fibre \(F_x\). Hence \(\mu=\nu\) exactly when the conditional distributions agree on every active fibre, proving the first assertion.

By definition of the conditional-uniform lift,
\[
(Q_m\pi)(a)=\frac{\pi(x)}{|F_x|}
\qquad(a\in F_x).
\]
Therefore \(\mu=Q_m\pi\) is equivalent to \(\mu_x=u_x\) on each active fibre. Since all spaces are finite, the elementary property of relative entropy gives
\[
D_{\mathrm{KL}}(\mu\|Q_m\pi)=0
\quad\Longleftrightarrow\quad
\mu=Q_m\pi,
\]
which proves the second assertion. Applying the same equivalence to \(\nu\), under the common pushforward \(
\pi=P_m\mu=P_m\nu\), gives the final statement.证毕。
\end{proof}

该量刻画“两条微态在何种分辨率上首次被折叠读出区分”。由此可把“类内区别的层级”理解为一个以 \(m\) 为刻度的分割深度：越大的 \(m_{\mathrm{sep}}\) 表示越晚（越深）才会从折叠读出中显现差别。该层级与前缀超度量一致地组织“柱拓扑—可判定性—分辨率门槛”的语义，而不需要额外的背景维度。

\subsubsection{序列层可逆化：纤维残差的时间化外置}\label{subsubsec:frt-sequence-inversion}
窗口层折叠 \(\Fold_m\) 典型为多对一，但其在序列层的滑动块码提升 \(\Phi_m\) 可在阈值后恢复可逆性：定理 \ref{thm:Phi_m-conjugacy-threshold} 证明，当 \(m\ge 3\) 时 \(\Phi_m\) 为拓扑共轭，且存在一个逆码 \(\Psi_m\) 使
\[
\Psi_m\circ\Phi_m=\mathrm{id},\qquad \Phi_m\circ\Psi_m=\mathrm{id},
\]
并且 \(\Psi_m\) 的记忆为 \(m-1\)。这给出“残差 \(\Rightarrow\) 时间深度”的严格译文：窗口层被折叠压缩掉的纤维自由度并未消失，而是被系统性迁移为序列层逆码所需的有限记忆深度。\par
在同一口径下，同步时间 \(T_m\) 以在线方式度量“候选集收缩到单点所需的最短读入深度”（定义 \ref{def:Ym-sync-time}），其意义可表述为“统计型 \(\mathsf{NULL}\) 首次消失”的时刻（见小节 \ref{subsec:Ym-zeta-spectrum} 的定义后解释段落）。因此，\(\Psi_m\) 的有限记忆与 \(T_m\) 的在线收缩深度共同提供了两种等价的“残差展开深度”刻度：前者是离线逆码口径，后者是在线滤波口径。

\subsubsection{最小在线延迟：残差外化的最小时间量子}\label{subsubsec:frt-min-online-delay}
若要求折叠在在线、因果与有限状态意义下实现，则残差的外化存在不可约的最小延迟。定理 \ref{thm:online-delay-fold} 给出一个延迟 \(3\) 的 \(\Fold_\infty\) 单趟在线实现；命题 \ref{prop:pom-delay-min} 进一步证明该常数是最小的，即
\[
\delta_{\min}(\Fold_\infty)=3.
\]
该结论把“时间的最小量子”钉成一个可审计的机制下界：延迟并非工程选择，而是由 Zeckendorf 规范化的跨位传播与同步状态强制产生。\par
并且该“时间缺陷”与“空间缺陷”（最大纤维多重度）互不替代：增加空间无法把 \(\delta_{\min}\) 压到 \(2\)，而压低延迟也不能消去纤维多重度（注 \ref{rem:pom-space-time-orth}）。在本文的统一口径下，二者与谱指纹共同组成“空间—时间—谱三指纹”（定义 \ref{def:pom-fingerprint}）。
\begin{remark}[时间并非附加几何维度，而是显影深度]\label{rem:frt-time-as-gap-depth}
本小节中的“时间”不应理解为外加的背景坐标，而应理解为纤维差异在协议中首次显影所需的解析深度。折叠分离深度 \(m_{\mathrm{sep}}\)、序列层逆码的有限记忆，以及
\[
\delta_{\min}(\Fold_\infty)=3
\]
所记录的是同一事实：差异并非不存在，而是必须推进到足够深的前缀分辨率、足够长的前视窗口或足够细的同步层级之后，才首次成为可审计对象。
\par
因此，后继轨道与极限对象之间的断层在时间化语义中并不对应一条额外连续轴，而对应一个不可压缩的最小显影深度。所谓残差外化的“时间读法”，正是对这一最小深度的协议化记录。
\end{remark}


\subsubsection{同步时间尾与谱证书}\label{subsubsec:frt-sync-tail-spectrum}
同步时间不仅可定义，还具有可审计的尾律：在最大熵测度下，同步尾满足指数界
\[
\nu_m(T_m>n)\le C_m e^{-\varepsilon_m n}
\]
（定理 \ref{thm:Ym-sync-tail}）。该结论将“残差寿命”对象化为一条时间统计，并把其衰减率与歧义壳的有限图结构锁定在同一谱对象上（定理 \ref{thm:Ym-sync-tail} 给出的矩阵幂表达与极限率 \(\varepsilon_m\)）。在本文的 Zeckendorf 折叠中，\(m\ge 3\) 时歧义壳为空（见小节 \ref{subsec:Ym-amb-transient} 的结论与备注），从而同步时间在有限处截断；该退化作为一致性检查被保留，强调“同步成本”在本模型中是有限可审计常数而非额外熵损失（同小节）。

\subsubsection{高阶碰撞矩与 \texorpdfstring{$\mathrm{MOM}_q$}{MOMq}：类内结构的不变量族}\label{subsubsec:frt-moments}
纤维多重度的全体分布 \(\{d_m(x)\}_{x\in X_m}\) 可由高阶碰撞矩统一编码：对整数 \(q\ge 2\)，定义
\[
S_q(m):=\sum_{x\in X_m}d_m(x)^q.
\]
其意义是“\(q\) 体仍被折叠为同一稳定类型”的重合计数。POM 语言将其提升为一个可编译门：\(\mathrm{MOM}_q\) 把底层在线核编译为 \(q\)-重 fiber-product moment-kernel，从而在有限状态上读出 \(S_q(m)\) 的递推与 Perron 主尺度（定理 \ref{thm:pom-collision-kernel-ak}；并见小节 \ref{subsec:pom-mom-rewrite-normal}）。因此，高阶矩并非启发式特征，而是纤维谱的可审计不变量族：它们在有限尺度上决定 \(\sum_x w_m(x)^q=S_q(m)/2^{qm}\) 的全部 Rényi 指纹率，并进一步在统一闭环中与 \(\zeta\)/primitive/谱对象对齐（定理 \ref{thm:pom-unified-closure}）。

\subsubsection{频域封装：\texorpdfstring{$\zeta$}{zeta} 接口与时间—空间交换图}\label{subsubsec:frt-zeta-packaging}
残差的“时间展开”不仅可在同步尾与在线延迟中被直接读出，还可在频域以 \(\zeta\) 形式被封装。命题 \ref{prop:time-space-commuting-diagram} 给出严格交换图 \(Z=E\circ M\circ T\)：时间侧以迹序列编码迭代统计，空间侧以 primitive 轨道谱实现 Euler 乘积与 Möbius 反演，从而把“时间读出”与“空间原子分解”统一到同一 \(\zeta\)-行列式接口。结合定理 \ref{thm:pom-unified-closure}，在线核的时间缺陷 \(\delta_{\min}\)、碰撞矩谱的 Perron 根族 \(\{r_q\}\) 与第 \ref{sec:zeta-finite-part} 节的谱/有限部指纹由同一可计算接口闭合。\par
在该闭合口径下，“时间作为残差”可被精确理解为：折叠所诱导的核方向信息不能在协议中被隐式抹除，只能在可审计的正交通道中外置（记忆深度、同步尾或额外寄存器轴），并以谱对象/递推/尾律等形式给出可复核证书。

\subsubsection{表观随机：冻结相与确定性粗糙统计}\label{subsubsec:frt-apparent-randomness-freezing-rough-statistics}

小节 \ref{subsubsec:groupoid-zeckendorf-multiscale-rough-readout} 已表明：在完全确定的隐藏动力学与残差相容的多尺度读出下，可见量已经可以连续而处处不可导。现转入统计层面，讨论这一类粗糙可见层在有限分辨率下为何会呈现出随机样输出。就本文既有对象而言，所需成分事实上已经齐备：一方面，最小在线延迟 \(\delta_{\min}(\Fold_\infty)=3\) 与同步尾指数界把残差的时间寿命对象化；另一方面，纤维多重度 \(d_m(x)\)、高阶碰撞矩 \(S_q(m)\)、冻结阈值与 escort 集中又把残差的空间谱对象化；再经 \(\zeta\)-接口、finite-part 与射影压力算子，这条链已在谱侧闭合。因此，本段不再把“像随机”当作修辞，而是把它改写为一个可审计判别：单点边缘不足以恢复机制，必须进入块碰撞、冻结相、mixed collision 与 finite-part 残差这些高阶不变量层。

\paragraph{块碰撞与最小判别}
为把粗糙可见量落到有限窗口协议上，对每个 \(m\ge 1\) 与 \(\omega\in\Omega_m\)，取任一延拓 \(\widehat\omega\in\Omega_\infty\)，并定义
\[
Y_{\omega,m}(t):=Y^{(\le m)}_{\widehat\omega}(t),
\]
其中 \(Y^{(\le m)}\) 沿用定义 \ref{def:groupoid-zeckendorf-multiscale-rough-readout} 的截断读出。由于该量只依赖于前 \(m\) 层残差标签，故与延拓的选择无关。固定采样时刻
\[
t_0,\dots,t_{\ell-1}
\]
及分箱映射
\[
Q_{m,\varepsilon}:\RR\to \mathcal A_{m,\varepsilon},
\]
定义长度为 \(\ell\) 的可见块读出
\[
\Xi^{(\ell)}_{m,\varepsilon}(\omega)
:=
\bigl(
Q_{m,\varepsilon}(Y_{\omega,m}(t_0)),
\dots,
Q_{m,\varepsilon}(Y_{\omega,m}(t_{\ell-1}))
\bigr)
\in \mathcal A_{m,\varepsilon}^{\ell}.
\]
在均匀微态基线下，相应块分布为
\[
w^{(\ell)}_{m,\varepsilon}(u)
:=
2^{-m}\#\bigl\{\omega\in\Omega_m:\ \Xi^{(\ell)}_{m,\varepsilon}(\omega)=u\bigr\},
\qquad
u\in \mathcal A_{m,\varepsilon}^{\ell}.
\]
当 \(\ell=1\) 且 \(Q_{m,\varepsilon}\) 退化为稳定类型读出时，上式正回到
\[
w_m(x)=\frac{d_m(x)}{2^m}.
\]
因此这一步并非引入新体系，而只是把既有的单点分布扩展为块口径。

\begin{definition}[块碰撞量与块缺口]\label{def:frt-block-collision-gap}
对整数 \(q\ge 2\)，定义块 \(q\)-重碰撞量
\[
\mathrm{Col}^{(\ell)}_q(m,\varepsilon)
:=
\sum_{u\in \mathcal A_{m,\varepsilon}^{\ell}}
\bigl(w^{(\ell)}_{m,\varepsilon}(u)\bigr)^q.
\]
并定义与同边缘 i.i.d.\ 基线的偏离量
\[
\Gamma^{(\ell)}_q(m,\varepsilon)
:=
\log \mathrm{Col}^{(\ell)}_q(m,\varepsilon)
-\ell\log \mathrm{Col}^{(1)}_q(m,\varepsilon).
\]
\end{definition}

\begin{proposition}[同边缘 i.i.d.\ 基线的乘法分解]\label{prop:frt-iid-block-factorization}
\leanverified{paper\_frt\_iid\_block\_factorization\_seeds}
\leanverified{paper\_frt\_iid\_block\_factorization\_package}
\leanverified{paper\_frt\_iid\_block\_factorization}
设
\[
\nu^{(\ell)}_{m,\varepsilon}
:=
\bigl(w^{(1)}_{m,\varepsilon}\bigr)^{\otimes \ell}
\]
为与 \(w^{(1)}_{m,\varepsilon}\) 具有相同单点边缘的 i.i.d.\ 参考块分布，则对任意整数 \(q\ge 2\)，有
\[
\sum_{u\in \mathcal A_{m,\varepsilon}^{\ell}}
\bigl(\nu^{(\ell)}_{m,\varepsilon}(u)\bigr)^q
=
\Bigl(\mathrm{Col}^{(1)}_q(m,\varepsilon)\Bigr)^\ell.
\]
因而
\[
\Gamma^{(\ell)}_q(m,\varepsilon)=0
\]
当且仅当块 \(q\)-碰撞层对同边缘 i.i.d.\ 基线完全因子化。
\end{proposition}
\begin{proof}
对 \(u=(u_0,\dots,u_{\ell-1})\in \mathcal A_{m,\varepsilon}^{\ell}\)，有
\[
\nu^{(\ell)}_{m,\varepsilon}(u)
=
\prod_{j=0}^{\ell-1} w^{(1)}_{m,\varepsilon}(u_j).
\]
因此
\[
\sum_{u\in \mathcal A_{m,\varepsilon}^{\ell}}
\bigl(\nu^{(\ell)}_{m,\varepsilon}(u)\bigr)^q
=
\sum_{u_0,\dots,u_{\ell-1}}
\prod_{j=0}^{\ell-1}\bigl(w^{(1)}_{m,\varepsilon}(u_j)\bigr)^q
=
\prod_{j=0}^{\ell-1}
\sum_{a\in \mathcal A_{m,\varepsilon}}
\bigl(w^{(1)}_{m,\varepsilon}(a)\bigr)^q,
\]
即得所示恒等式。证毕。
\end{proof}

\begin{remark}[表观随机的最小判别]\label{rem:frt-block-collision-gap-meaning}
命题 \ref{prop:frt-iid-block-factorization} 说明：单点边缘可以与某个独立噪声基线完全一致，而块碰撞缺口 \(\Gamma_q^{(\ell)}\) 一旦非零，就已足以排除“同边缘 i.i.d.”这一解释。对 Fold 诱导的粗糙可见层，这种块层非因子化并非附会，而是由最小在线延迟 \(\delta_{\min}(\Fold_\infty)=3\)（定理 \ref{thm:online-delay-fold} 与命题 \ref{prop:pom-delay-min}）以及同步尾所对应的有限记忆链自然预示的。换言之，在本文口径下，“看起来像随机”与“真正无记忆”本来就是两件不同的事情。
\end{remark}

\paragraph{冻结相与连续压力包络}
单点层的高阶碰撞已由
\[
S_q(m):=\sum_{x\in X_m}d_m(x)^q
\]
编码；其归一化碰撞量为
\[
\sum_{x\in X_m} w_m(x)^q
=
\frac{S_q(m)}{2^{qm}}.
\]
为避免与附录 \ref{subsubsec:fold-fiber-groundstate-freezing-oracle-capacity} 中的未归一化压力记号冲突，下文记
\[
\tau(q):=\limsup_{m\to\infty}\frac{1}{m}\log S_q(m),
\qquad
P(q):=\tau(q)-q\log 2.
\]
若 Perron 根 \(r_q\) 存在，则 \(\tau(q)=\log r_q\)，故
\[
P(q)=\log r_q-q\log 2.
\]
这正是“单点 \(q\)-碰撞概率的指数尺度”。

\begin{theorem}[冻结相与 escort 集中]\label{thm:frt-apparent-randomness-freezing-escort}
设附录 \ref{subsubsec:fold-fiber-groundstate-freezing-oracle-capacity} 中的严格基态间隙条件
\[
\alpha_*^{(2)}<\alpha_*
\]
成立，并记冻结阈值
\[
a_{\mathrm c}:=\frac{\log\varphi-g_*}{\alpha_*-\alpha_*^{(2)}}.
\]
则对每个整数 \(q>a_{\mathrm c}\)，都有
\[
\tau(q)=g_*+q\alpha_*,
\qquad
P(q)=g_*+q(\alpha_*-\log 2).
\]
此外，对 escort 分布
\[
\pi_{m,q}(x):=\frac{d_m(x)^q}{S_q(m)}
\]
及基态集合
\[
A_m^{\max}:=\{x\in X_m:\ d_m(x)=M_m\},
\qquad
M_m:=\max_{x\in X_m}d_m(x),
\]
存在常数 \(\Delta(q)>0\) 使
\[
\pi_{m,q}(A_m^{\max})
\ge
1-\exp\bigl(-m\Delta(q)+o(m)\bigr).
\]
\end{theorem}
\begin{proof}
第一式正是定理 \ref{thm:fold-pressure-freezing-threshold} 的内容；第二式由 \(P(q)=\tau(q)-q\log 2\) 直接得到。escort 的指数集中则是定理 \ref{thm:fold-escort-groundstate-concentration} 的重写。证毕。
\end{proof}
\leanverified{paper\_frt\_apparent\_randomness\_freezing\_escort}

\begin{remark}[连续射影压力不是事后插值]\label{rem:frt-apparent-randomness-projective-pressure}
附录 \ref{app:pom-projective-pressure-operator} 已用定义 \ref{def:pom-projective-transfer-operator} 的射影转移算子 \(\mathcal L_q\) 给出一条实参数压力曲线。记其 Perron 主特征值为 \(\lambda(q)\)，并沿注 \ref{rem:pom-projective-pressure-route} 定义
\[
\Lambda(t):=\log \lambda(t+1)-\log |A_{\mathrm{out}}|.
\]
则由定理 \ref{thm:pom-projective-operator-degenerates-to-moment-kernel}，在整数点 \(t=q-1\) 上有
\[
\Lambda(q-1)=\log r_q-\log 2=P(q)+(q-1)\log 2.
\]
因此，连续压力并不是对离散矩谱做任意平滑，而是对整数锚点 \(\log r_q\) 的算子实现；归一化碰撞压力 \(P(q)\) 只是该连续曲线的一个仿射重标。对粗糙可见统计而言，离散冻结相、连续压力包络与 Legendre 型率函数因此处于同一条可审计链上。
\end{remark}

\paragraph{mixed collision 与 finite-part 残差}
若只看单一机制的 \(\{S_q(m)\}\)，仍不足以回答“两个都像随机的机制如何区分”。这时需要跨机制比较量。小节 \ref{subsubsec:pom-hoelder-mixed-kl-bridge} 已定义 mixed collision
\[
\mathrm{Col}^{(0,1)}_q(m)
:=
\sum_{x\in X_m}
\bigl(w_m^{(0)}(x)\bigr)^q
\bigl(w_m^{(1)}(x)\bigr)^q
\]
并证明其对固定 \(q\) 可由有限状态在线核编译为有限维非负矩阵幂读出（命题 \ref{prop:pom-mixed-collision-kernel-computable}）。另一方面，小节 \ref{subsec:pom-mom-rewrite-normal} 的定义 \ref{def:pom-moment-anom-signature} 与命题 \ref{prop:pom-em-mom-residual-swap} 又把 \(E_m\) 与 \(\mathrm{MOM}_q\) 的可观测不交换性压缩为 finite-part 标量残差
\[
\mathsf{Anom}^{(q)}_{\mathrm{mom}}(K;\theta).
\]
因此，“确定性粗糙”与“真噪声”的区分至少应同时考察块层非因子化、跨机制 overlap 与 finite-part 残差三条证书。

\begin{definition}[归一化 mixed collision 相似度]\label{def:frt-normalized-mixed-collision-similarity}
对任意整数 \(q\ge 1\)，定义
\[
\Xi_q^{(0,1)}(m)
:=
\frac{\mathrm{Col}^{(0,1)}_q(m)}
{\sqrt{\mathrm{Col}^{(0,0)}_q(m)\,\mathrm{Col}^{(1,1)}_q(m)}}.
\]
\end{definition}

\begin{proposition}[mixed collision 的 Cauchy 缺口]\label{prop:frt-mixed-collision-cauchy-gap}
对任意两分布 \(w_m^{(0)},w_m^{(1)}\in\Delta(X_m)\) 与任意整数 \(q\ge 1\)，有
\[
0\le \Xi_q^{(0,1)}(m)\le 1.
\]
并且等号 \(\Xi_q^{(0,1)}(m)=1\) 成立当且仅当
\[
w_m^{(0)}=w_m^{(1)}
\qquad\textup{于 }X_m\textup{ 上逐点相等}.
\]
\end{proposition}
\begin{proof}
令
\[
a_x:=\bigl(w_m^{(0)}(x)\bigr)^q,
\qquad
b_x:=\bigl(w_m^{(1)}(x)\bigr)^q.
\]
则
\[
\mathrm{Col}^{(0,1)}_q(m)=\sum_{x\in X_m}a_xb_x,
\qquad
\mathrm{Col}^{(0,0)}_q(m)=\sum_{x\in X_m}a_x^2,
\qquad
\mathrm{Col}^{(1,1)}_q(m)=\sum_{x\in X_m}b_x^2.
\]
由 Cauchy--Schwarz 不等式，
\[
\Bigl(\sum_{x\in X_m}a_xb_x\Bigr)^2
\le
\Bigl(\sum_{x\in X_m}a_x^2\Bigr)
\Bigl(\sum_{x\in X_m}b_x^2\Bigr),
\]
故 \(0\le \Xi_q^{(0,1)}(m)\le 1\)。若等号成立，则 \((a_x)\) 与 \((b_x)\) 必线性相关；再利用两者均来自概率分布，可知比例常数只能为 \(1\)，从而 \(w_m^{(0)}=w_m^{(1)}\)。反向显然。证毕。
\end{proof}
\leanverified{paper\_frt\_mixed\_collision\_cauchy\_gap}

\begin{remark}[三类互补证书]\label{rem:frt-apparent-randomness-three-certificates}
在本文语境下，\(\Gamma_q^{(\ell)}(m,\varepsilon)\) 检测的是块层是否真正像同边缘 i.i.d.\ 噪声；\(\Xi_q^{(0,1)}(m)\) 检测的是两条机制在 \(q\)-阶尾部几何上有多相似；\(\mathsf{Anom}^{(q)}_{\mathrm{mom}}(K;\theta)\) 检测的是“先条件期望后取矩”与“先进入 moment 世界再平均”是否可交换。三者都通过时，粗糙确定性系统才可能在可审计意义下与某个随机基线难以区分；而任一层的严格缺口都已足以把“表观随机”与“本体噪声”分开。
\end{remark}

\paragraph{window-\texorpdfstring{$6$}{6} 锚点}
为避免上述对象显得过于抽象，回到本文最重要的固定窗锚点 \(m=6\)。推论 \ref{cor:fold-bin-m6-fiber-hist} 已给出二进制区间折叠的刚性直方图
\[
2:8,\qquad 3:4,\qquad 4:9.
\]
这意味着：window-\(6\) 的 \(21\) 个可见稳定类型并不等势，因而其统计外观不可能来自任意自由群作用或线性陪集划分，而必然携带真实的纤维层级。

\begin{proposition}[window-\texorpdfstring{$6$}{6} 的固定窗冻结律]\label{prop:frt-window6-fixed-freezing-law}
\leanverified{paper\_frt\_window6\_fixed\_freezing\_law}
对 window-\(6\) 的 bin-fold，定义
\[
S_q^{(6)}:=\sum_{x\in X_6}\bigl(d_6^{\mathrm{bin}}(x)\bigr)^q.
\]
则对任意整数 \(q\ge 1\) 有
\[
S_q^{(6)}=8\cdot 2^q+4\cdot 3^q+9\cdot 4^q.
\]
因此其固定窗未归一化压力
\[
\tau_q^{(6)}:=\log S_q^{(6)}
\]
满足精确分解
\[
\tau_q^{(6)}
=
q\log 4+\log 9
+\log\!\left(
1+\frac{4}{9}\Bigl(\frac{3}{4}\Bigr)^q+\frac{8}{9}\,2^{-q}
\right).
\]
特别地，当 \(q\to\infty\) 时，
\[
\tau_q^{(6)}
=
q\log 4+\log 9
+\frac{4}{9}\Bigl(\frac{3}{4}\Bigr)^q
+O\!\left(\Bigl(\frac{9}{16}\Bigr)^q\right),
\]
并且相应 escort 质量在最大纤维层 \(d=4\) 上收敛到 \(1\)。
\end{proposition}
\begin{proof}
由推论 \ref{cor:fold-bin-m6-fiber-hist} 的直方图直接得
\[
S_q^{(6)}=8\cdot 2^q+4\cdot 3^q+9\cdot 4^q.
\]
提出 \(9\cdot 4^q\) 即得
\[
S_q^{(6)}
=
9\cdot 4^q
\left(
1+\frac{4}{9}\Bigl(\frac{3}{4}\Bigr)^q+\frac{8}{9}\,2^{-q}
\right),
\]
从而得到 \(\tau_q^{(6)}\) 的精确分解。再由 \(\log(1+t)=t+O(t^2)\) 以及
\[
\Bigl(\frac{3}{4}\Bigr)^{2q},\ 2^{-q}
=
O\!\left(\Bigl(\frac{9}{16}\Bigr)^q\right),
\]
得到渐近式。最后，escort 在层 \(d=4\) 上的质量为
\[
\frac{9\cdot 4^q}{8\cdot 2^q+4\cdot 3^q+9\cdot 4^q},
\]
当 \(q\to\infty\) 时显然趋于 \(1\)。证毕。
\end{proof}

\begin{corollary}[window-\texorpdfstring{$6$}{6} 的二点碰撞常数]\label{cor:frt-window6-second-collision-constant}
在命题 \ref{prop:frt-window6-fixed-freezing-law} 的设定下，
\[
S_2^{(6)}=8\cdot 4+4\cdot 9+9\cdot 16=212,
\]
从而
\[
\mathrm{Col}^{(6)}_2
:=
\sum_{x\in X_6}\left(\frac{d_6^{\mathrm{bin}}(x)}{2^6}\right)^2
=
\frac{S_2^{(6)}}{2^{12}}
=
\frac{53}{1024}.
\]
\leanverified{paper\_frt\_window6\_second\_collision\_constant}
\end{corollary}

\begin{proposition}[冻结 escort 强制 sticky 或基态 carrier 增殖]\label{distill:wang_zahl-sticky_reduction_and_scale_saturated_closure-freezing-forces-sticky-or-groundstate-proliferation}
设定理 \ref{thm:frt-apparent-randomness-freezing-escort} 的严格基态间隙成立，并取整数 \(q>a_{\mathrm c}\)。记
\[
\pi_{m,q}(x)=\frac{d_m(x)^q}{S_q(m)},
\qquad
A_m^{\max}:=\{x\in X_m:d_m(x)=M_m\}.
\]
于是存在 \(\Delta(q)>0\) 与 \(\varepsilon_m=\exp(-m\Delta(q)+o(m))\)，使
\[
\pi_{m,q}(A_m^{\max})\ge1-\varepsilon_m.
\]
令 \(\gamma_{m,h}:X_m\to Y_h\) 为任意有限 coarse carrier 映射，并设
\[
G_{m,h}^{\max}:=\gamma_{m,h}(A_m^{\max}).
\]
则
\[
\max_{y\in Y_h}\pi_{m,q}(\gamma_{m,h}^{-1}(y))
\ge
\frac{1-\varepsilon_m}{|G_{m,h}^{\max}|}.
\]
因此，对任意整数 \(s\ge1\)，若
\[
|G_{m,h}^{\max}|<(1-\varepsilon_m)2^s,
\]
则 \(X_m\) 在 \(\gamma_{m,h}\) 上为 \(s\)-sticky。反过来，若 \(X_m\) 在同一 carrier 映射上为 \(s\)-non-sticky，则必有
\[
|G_{m,h}^{\max}|\ge(1-\varepsilon_m)2^s.
\]
也就是说，冻结相中的 non-sticky 分支不能由基态质量弥散解释；它强制近全部 escort 质量所在的最大纤维层在 coarse carrier 上发生指数级增殖。
\end{proposition}
\begin{proof}
由定理 \ref{thm:frt-apparent-randomness-freezing-escort}，存在所述 \(\varepsilon_m\) 使
\[
\pi_{m,q}(A_m^{\max})\ge1-\varepsilon_m.
\]
集合 \(A_m^{\max}\) 被 carrier 映射 \(\gamma_{m,h}\) 分解到 \(G_{m,h}^{\max}\) 个 occupied carrier 中。因此
\[
1-\varepsilon_m
\le
\pi_{m,q}(A_m^{\max})
\le
\sum_{y\in G_{m,h}^{\max}}
\pi_{m,q}(\gamma_{m,h}^{-1}(y))
\le
|G_{m,h}^{\max}|
\max_{y\in Y_h}\pi_{m,q}(\gamma_{m,h}^{-1}(y)).
\]
移项即得最大 carrier 质量下界。

若 \(|G_{m,h}^{\max}|<(1-\varepsilon_m)2^s\)，则上式给出某个 carrier 的 \(\pi_{m,q}\)-质量严格大于 \(2^{-s}\)。当 \(E=X_m\) 时，定义 \ref{def:distill-folding-integer-depth-sticky-test} 中的 escort 质量正是 \(\pi_{m,q}\)，故 \(X_m\) 为 \(s\)-sticky。其逆否命题即为 non-sticky 时的基态 carrier 增殖下界。证毕。
\end{proof}

\begin{proposition}[冻结非 sticky 分支的有限分辨率指数预算下界]\label{distill:wang_zahl-worst_counterexample_exponent_bootstrapping-frt-nonsticky-carrier-exponent-budget}
设定理 \ref{thm:frt-apparent-randomness-freezing-escort} 的严格基态间隙成立，并取整数 \(q>a_{\mathrm c}\)。沿用
\[
\pi_{m,q}(x)=\frac{d_m(x)^q}{S_q(m)},
\qquad
A_m^{\max}=\{x\in X_m:d_m(x)=M_m\},
\qquad
\varepsilon_m=\exp(-m\Delta(q)+o(m))
\]
的记号，并取 \(m\) 充分大使 \(0<\varepsilon_m<1\)。令
\[
\gamma_{m,h}:X_m\to Y_h
\]
为任意有限 coarse carrier 映射，且
\[
G_{m,h}^{\max}:=\gamma_{m,h}(A_m^{\max}).
\]
给定任意尺度归一化常数 \(L_h>0\)，定义 ground-state carrier 指数
\[
D_{m,h}^{\max}:=\frac{1}{L_h}\log_2 |G_{m,h}^{\max}|
\]
以及有限分辨率断言
\[
\mathsf K_{m,h}(D):\quad D_{m,h}^{\max}\le D.
\]
对整数 \(s\ge1\)，记 \(\mathsf{NST}_{m,h}(s)\) 表示：在定义 \ref{def:distill-folding-integer-depth-sticky-test} 的 sticky 测试中，以 \(E=X_m\)、carrier 映射 \(\gamma_{m,h}\) 与 escort 权重 \(\pi_{m,q}\) 为输入时，\(X_m\) 为 \(s\)-non-sticky。

则有严格预算下界
\[
\mathsf{NST}_{m,h}(s)
\quad\Longrightarrow\quad
|G_{m,h}^{\max}|\ge (1-\varepsilon_m)2^s,
\]
从而
\[
\mathsf{NST}_{m,h}(s)
\quad\Longrightarrow\quad
D_{m,h}^{\max}\ge
\frac{s+\log_2(1-\varepsilon_m)}{L_h}.
\]
特别地，若 \(D\ge0\)、\(\alpha>0\)、\(\eta>0\)，且
\[
s\ge (D+\alpha+\eta)L_h+\log_2\!\bigl((1-\varepsilon_m)^{-1}\bigr),
\]
则任何 \(s\)-non-sticky 分支都强制
\[
D_{m,h}^{\max}\ge D+\alpha+\eta,
\]
因而不满足 \(\mathsf K_{m,h}(D+\alpha)\)。等价地，若 \(\mathsf K_{m,h}(D_0)\) 成立，则对一切满足
\[
s>D_0L_h+\log_2\!\bigl((1-\varepsilon_m)^{-1}\bigr)
\]
的整数 \(s\)，该有限分辨率对象必为 \(s\)-sticky。
\end{proposition}

\begin{proposition}[表观随机审计的有限测试稳定账本]\label{prop:distill-euclid-elements-frt-finite-test-stability-ledger}
\textbf{Dependency status: conditional audit.}
取有限测试族
\[
\mathcal T_{\rm blk}\subset
\{(m,\varepsilon,\ell,q):m\ge1,\ \ell\ge1,\ q\ge2\}
\]
与
\[
\mathcal T_{\rm mix}\subset
\{(m,q):m\ge1,\ q\ge1\},
\]
其中所有出现的 \(\Gamma_q^{(\ell)}(m,\varepsilon)\) 与 \(\Xi_q^{(0,1)}(m)\) 均按定义 \ref{def:frt-block-collision-gap} 与定义 \ref{def:frt-normalized-mixed-collision-similarity} 解释。定义该有限测试账本的残差预算为
\[
\mathfrak R(\mathcal T_{\rm blk},\mathcal T_{\rm mix})
:=\max\left(
\{0\}\cup
\{|\Gamma_q^{(\ell)}(m,\varepsilon)|:(m,\varepsilon,\ell,q)\in\mathcal T_{\rm blk}\}
\cup
\{1-\Xi_q^{(0,1)}(m):(m,q)\in\mathcal T_{\rm mix}\}
\right).
\]
则：
\begin{enumerate}
  \item 若 \(\mathcal T_{\rm blk}\subseteq\mathcal T'_{\rm blk}\) 且 \(\mathcal T_{\rm mix}\subseteq\mathcal T'_{\rm mix}\)，则
  \[
  \mathfrak R(\mathcal T_{\rm blk},\mathcal T_{\rm mix})
  \le
  \mathfrak R(\mathcal T'_{\rm blk},\mathcal T'_{\rm mix});
  \]
  \item 对任意有理容忍阈值 \(\eta\in\QQ_{\ge0}\)，该有限测试族全部通过
  \[
  |\Gamma_q^{(\ell)}(m,\varepsilon)|\le\eta,
  \qquad
  1-\Xi_q^{(0,1)}(m)\le\eta
  \]
  当且仅当
  \[
  \mathfrak R(\mathcal T_{\rm blk},\mathcal T_{\rm mix})\le\eta;
  \]
  \item 若 \((\mathcal T^{(R)}_{\rm blk},\mathcal T^{(R)}_{\rm mix})_{R\ge1}\) 是递增的有限测试账本列，则
  \[
  R\longmapsto
  \mathfrak R(\mathcal T^{(R)}_{\rm blk},\mathcal T^{(R)}_{\rm mix})
  \]
  单调不降。若该数列对所有 \(R\) 均不超过 \(\eta\)，则并集中的每一个已声明有限测试均以阈值 \(\eta\) 通过。
\end{enumerate}
该结论只是有限测试稳定账本；它不声称未声明的块长度、分箱或机制对已经满足全局独立性或完整随机基线。
\end{proposition}

\begin{proposition}[有限 packet 平均的确定性冻结证书]\label{prop:distill-bourgain-frt-finite-packet-freezing-extraction}
令 \(\mathcal P\) 为非空有限 packet 族，令 \(\alpha_P\in\QQ_{\ge0}\) 满足 \(\sum_{P\in\mathcal P}\alpha_P=1\)。对每个 \(P\in\mathcal P\)，给定一个已离线计算的残差值 \(R(P)\in\QQ_{\ge0}\) 与一个生存谓词 \(\chi(P)\in\{0,1\}\)。给定 \(\theta,\eta\in\QQ_{>0}\)。若有限审计账本证明
\[
\sum_{P:\chi(P)=1}\alpha_P\ge\theta,
\qquad
\sum_{P:\chi(P)=1}\alpha_P R(P)\le\theta\eta,
\]
则存在 \(P_\ast\in\mathcal P\) 使
\[
\alpha_{P_\ast}>0,\qquad \chi(P_\ast)=1,\qquad R(P_\ast)\le\eta.
\]
因此，任何以平均、随机抽样或大偏差语言得到的有限 packet 结论，只有在选定并存储这样的 \(P_\ast\) 后，才成为本文意义下的确定性粗糙统计证书。若没有给出上述 \(\theta\)-质量账本或没有冻结 \(P_\ast\)，则至多说明当前有限 packet 测试尚未完成证书抽取，而不能声称所有未测试 packet 或其它机制已经通过随机性判别。
\end{proposition}
\begin{proof}
设反面成立，即每个满足 \(\alpha_P>0\) 且 \(\chi(P)=1\) 的 packet 都有 \(R(P)>\eta\)。由于所有数据为有理数且 \(\mathcal P\) 有限，求和可逐项比较，得到
\[
\sum_{P:\chi(P)=1}\alpha_P R(P)
>
\eta\sum_{P:\chi(P)=1}\alpha_P
\ge
\eta\theta,
\]
这与第二个审计不等式 \(\sum_{\chi(P)=1}\alpha_P R(P)\le\theta\eta\) 矛盾。因此存在正权且生存的 \(P_\ast\) 满足 \(R(P_\ast)\le\eta\)。

该证明只使用有限求和与显式生存谓词；它没有使用 packet 族外的任何概率模型。故冻结后的 \((P_\ast,R(P_\ast),\chi(P_\ast))\) 是确定性证书，而缺少冻结步骤时，平均界本身仍只是有限族内部的存在性信息。证毕。
\end{proof}

\begin{proof}
第一条由最大值取在更大集合上立得。第二条是最大值定义的直接等价：最大残差不超过 \(\eta\) 当且仅当每个被列入账本的残差均不超过 \(\eta\)。第三条由第一条应用于递增测试族得到单调性；若所有有限阶段的最大残差都不超过 \(\eta\)，任取并集中的一个测试，它属于某个有限阶段，故在该阶段已经通过。证毕。
\end{proof}

\begin{proof}
由命题 \ref{distill:wang_zahl-sticky_reduction_and_scale_saturated_closure-freezing-forces-sticky-or-groundstate-proliferation} 的逆否结论，若同一 carrier 映射 \(\gamma_{m,h}\) 上的 \(X_m\) 为 \(s\)-non-sticky，则最大纤维层在 coarse carrier 上必须满足
\[
|G_{m,h}^{\max}|\ge (1-\varepsilon_m)2^s.
\]
这正是第一式。两边取 \(\log_2\) 并除以 \(L_h\)，得到
\[
D_{m,h}^{\max}
=\frac{1}{L_h}\log_2 |G_{m,h}^{\max}|
\ge
\frac{s+\log_2(1-\varepsilon_m)}{L_h}.
\]
若
\[
s\ge (D+\alpha+\eta)L_h+\log_2\!\bigl((1-\varepsilon_m)^{-1}\bigr),
\]
则上式右端至少为 \(D+\alpha+\eta\)，故
\[
D_{m,h}^{\max}\ge D+\alpha+\eta>D+\alpha.
\]
于是 \(\mathsf K_{m,h}(D+\alpha)\) 失败。

最后，若 \(\mathsf K_{m,h}(D_0)\) 成立而某个
\[
s>D_0L_h+\log_2\!\bigl((1-\varepsilon_m)^{-1}\bigr)
\]
仍落入 non-sticky 分支，则前述下界给出
\[
D_{m,h}^{\max}>
D_0,
\]
与 \(\mathsf K_{m,h}(D_0)\) 矛盾。因此该深度下只能落入定义 \ref{def:distill-folding-integer-depth-sticky-test} 的 sticky 分支。

该断言没有把朴素截断量当作不变量：\(\mathsf K_{m,h}(D)\) 只作用于冻结 escort 已经选出的 fold-aware 最大纤维层 \(A_m^{\max}\) 及其 coarse carrier 像 \(G_{m,h}^{\max}\)。因此，在本有限分辨率口径下，最坏反例若拒绝 sticky 闭合，就必须支付明确的 carrier 指数预算；这正是此处可严格证明的 worst-counterexample exponent bootstrapping 形式。证毕。
\end{proof}

\begin{proof}
直接把 \(q=2\) 代入命题 \ref{prop:frt-window6-fixed-freezing-law} 即得。证毕。
\end{proof}

\begin{remark}[本段的统计含义]\label{rem:frt-apparent-randomness-summary}
至此，粗糙可见层的“表观随机”已可压缩为一条闭链：低阶上一体分布可能近似扩散，从而在观测上显得随机；但一旦进入块碰撞、escort、mixed collision 与 finite-part 层，就会显露出由纤维谱、同步记忆与矩编译不交换性共同控制的结构残差。因而，在 Fold--fiber--time--spectrum 的统一语言里，更准确的说法不是“系统随机”，而是：系统的可见统计在有限分辨率下进入了一个由残差层级所支配的表观随机区。
\end{remark}



\subsection{三条定理链：从协议到谱指纹的可审计导航}\label{subsec:three-theorem-chains-rewrite-skeleton}

本小节将全文的关键结果按依赖关系压缩为三条可复核的定理链。其目的不是引入新的“总定理”，而是给出一条便于审计的阅读路径：读者可沿每条链逐项回到原章节的编号结论完成复核。

\paragraph{链 I：协议刚性链（设计空间 \texorpdfstring{$\Rightarrow$}{=>} 黄金端点 \texorpdfstring{$\Rightarrow$}{=>} 折叠接口）}
\begin{enumerate}
  \item \textbf{二元读出与时间刻度的协议固定：}扫描—投影生成器与前缀柱拓扑（第 \ref{sec:spg} 节，段落 \ref{par:spg-prefix-ultrametric}）。
  \item \textbf{折叠构造范式：}Fibonacci 加权值与 Zeckendorf 规范形定义二元折叠 \(\Fold_m\)（定义 \ref{def:fold-word}；命题 \ref{prop:fold-basic}）。
  \item \textbf{一般设计空间：}一般 Ostrowski 合法集合与一般折叠族 \(\Fold_m^{(\alpha)}\)（小节 \ref{subsec:ostrowski-generalized-fold}；定义 \ref{def:Xm-alpha}、\ref{def:fold-alpha}）。
  \item \textbf{二元退化：}二元字母表把参数轴退化到 \(\alpha=\varphi^{-1}\)，并把合法语法退化到禁词 \((11)\) 的黄金均值语法（命题 \ref{prop:Xm-alpha-cardinality}；小节 \ref{subsec:folding-fibonacci-stable-syntax}）。
  \item \textbf{有限窗口完备：}稳定语法的前缀值域双射证书 \(V_m:X_m\simeq\{0,\dots,F_{m+2}-1\}\)（引理 \ref{lem:zeckendorf-range-bijection}；第 \ref{sec:emergent-arithmetic} 节的值—地址同构接口）。
  \item \textbf{跨分辨率相容与自然性：}一般折叠中“先规范形再截断”的函子性与其在二元折叠中的唯一修正（定理 \ref{thm:normalize-then-truncate-functorial}；命题 \ref{prop:fold-omega-commute}）。
  \item \textbf{最坏情形强制非局域：}规范差 \(G_m\) 的最坏端满长翻转与尾补丁不完备性（定理 \ref{thm:fold-gauge-anomaly-max}；推论 \ref{cor:fold-tail-patch-incomplete}）。
  \item \textbf{对称性审计：}window-\(6\) 锚点下几何稳定子的 \(\ZZ_2\) 刚性与掩码 \(\pm 34\) 指纹（定理 \ref{thm:foldbin6-geo-stabilizer-z2}；推论 \ref{cor:foldbin6-geo-mask-34}）。
\end{enumerate}

\paragraph{链 II：残差即时间链（商结构 \texorpdfstring{$\Rightarrow$}{=>} 纤维残差 \texorpdfstring{$\Rightarrow$}{=>} 在线/同步深度）}
\begin{enumerate}
  \item \textbf{折叠商结构与纤维：}纤维多重度 \(d_m(x)=|\Fold_m^{-1}(x)|\) 与稳定压缩率（小节 \ref{subsec:folding-stable-compression-resolution}）。
  \item \textbf{残差的严格对象化：}信息账本恒等式与纤维 KL 分解（定理 \ref{thm:pom-kl-ledger}；引理 \ref{lem:pom-fiberwise-kl}）。
  \item \textbf{分辨率分离与前缀几何：}前缀超度量与“球即柱集”的可判定口径（段落 \ref{par:spg-prefix-ultrametric}；定义 \ref{def:xi-dyadic-prefix-depth-ultrametric}）。
  \item \textbf{序列层可逆化：}滑动块因子 \(\Phi_m\) 的共轭阈值与有限记忆逆码（定理 \ref{thm:Phi_m-conjugacy-threshold}）。
  \item \textbf{最小在线延迟：}单趟在线延迟 \(3\) 的实现与最小性（定理 \ref{thm:online-delay-fold}；命题 \ref{prop:pom-delay-min}）。
  \item \textbf{同步时间与残差寿命：}同步时间 \(T_m\) 的定义与同步尾指数预算（定义 \ref{def:Ym-sync-time}；定理 \ref{thm:Ym-sync-tail}）。
  \item \textbf{空间/时间缺陷正交：}最大纤维多重度与最小延迟分别给出空间与时间指纹，互不替代（注 \ref{rem:pom-space-time-orth}；定义 \ref{def:pom-fingerprint}）。
\end{enumerate}

\paragraph{链 III：不变量与谱链（高阶矩 \texorpdfstring{$\Rightarrow$}{=>} MOM 编译 \texorpdfstring{$\Rightarrow$}{=>} \texorpdfstring{$\zeta$}{zeta}/primitive 接口）}
\begin{enumerate}
  \item \textbf{高阶碰撞矩：}纤维多重度谱由 \(S_q(m)=\sum_x d_m(x)^q\) 统一编码（例如定义 \ref{def:pom-s5} 及其低阶同型定义）。
  \item \textbf{moment-kernel 编译：}对任意 \(q\ge 2\)，存在有限状态 moment-kernel 读出 \(S_q\) 的递推与 Perron 主尺度（定理 \ref{thm:pom-collision-kernel-ak}）。
  \item \textbf{共振与实现维数：}Hankel 秩坍缩提供共振机制的可审计判别（定理 \ref{thm:pom-moment-resonance}；并见定理 \ref{thm:pom-unified-closure} 的链式收束）。
  \item \textbf{时间—空间交换图：}迹序列与 primitive 轨道谱在 \(\zeta\) 行列式口径下的严格对齐（命题 \ref{prop:time-space-commuting-diagram}）。
  \item \textbf{谱/有限部桥接：}由机制核的 \(\zeta\) 与有限部常数抽取谱指纹，并与群扩张/扭曲通道闭合（第 \ref{sec:zeta-finite-part} 节）。
  \item \textbf{统一闭环：}在线核（时间指纹）—碰撞矩核（空间谱）—\(\zeta\)/primitive（谱封装）—几何维数谱在同一接口下闭合（定理 \ref{thm:pom-unified-closure}）。
\end{enumerate}



